\section{讨论}
\subsection{选择性约束与动力学预测的关系}
初始两步法通过反应度域的指纹估计减少动态拟合自由度，为稀疏酯化曲线提供了参数约束 \cite{nishimura2026}。本文的组分残差、消融与模拟结果共同表明，指纹信息的作用取决于组成路径模型是否充分。Bornadel 数据中，固定不含伪持续性的指纹选择性会保留 M 与 D 分布的系统偏差；稀疏 Åkerman 数据中，完全取消指纹约束则增加预测误差。PPK-FIP 用有限宽度先验兼顾选择性稳定性与路径修正，主要贡献是改善已测试案例的留出时间点预测。

指纹先验由同一训练集估计，因此 PPK-FIP 更适合解释为数据自适应的正则化估计，而非利用独立先验信息的完整贝叶斯推断。留一时间点验证评估了估计流程的预测表现，但尚未检验后验区间的覆盖率。BIC 权重也采用近似模型权重，不应直接解释为机理真实概率。近期混合神经动力学研究同样强调参数约束和可辨识性分析的重要性 \cite{giampiccolo2024,derooij2025}；本文将参数约束具体落实为 TMP 酯化指纹估计与显式反应路径。

\subsection{伪持续性与失活参数的解释范围}
Bornadel 数据中，伪持续性修正同时改善反应度域组成路径、时间域残差和留出预测，与 Mitchell 等对传质受限条件下表观选择性的解释一致 \cite{mitchell2023}。已知真值模拟进一步显示，省略伪持续性可造成选择性偏差，适度放宽指纹先验能够减小偏差。跨水活度迁移则支持 Tao 数据中组成路径修正的可迁移性。浓度轨迹仍属于间接证据，无法唯一排除其他未建模反应或传质过程。

伪持续性概率应视为给定酶、载体与传质条件下的有效参数。Sengupta 等在另一固定化脂肪酶酯化体系中揭示了吸附和孔结构对反应过程的影响 \cite{sengupta2026}，进一步说明载体条件变化可能影响表观动力学。将 Bornadel 估计的选择性或伪持续性概率用于不同搅拌速度、颗粒尺度或载体时，需要重新验证组成路径。

批内失活主要改善 Åkerman 100\,°C 的全数据拟合，但失活项未降低合并交叉验证误差，去除失活还略微降低了合并外推误差。浓度曲线对失活的支持取决于采样窗口是否覆盖速率衰减阶段。独立酶活测量、延长时间采样及重复温度条件可帮助区分酶失活与其他减速来源。Ribeiro 等的固定化脂肪酶热失活实验支持在合适体系中考虑一阶衰减 \cite{ribeiro2025}，不能代替当前 TMP 体系的失活验证。

\subsection{参数辨识与验证设计}
初始模型核对显示，部分文献报告值与统一实现之间存在差异，多起点优化也改变了部分模型的目标值。完整集总常数模型的平坦剖面说明，在当前数据范围内，单个速率参数可能缺乏独立约束。低拟合误差、稳定参数解释和准确预测是不同评价目标，模型选择需要分别检查优化敏感性、参数区间与留出误差。

本文的留一时间点验证主要检验已观测条件和时间范围附近的预测能力。前半段训练的外推误差明显增加，且 PPK-FIP 相对于 ORIG 的合并优势不再显著，说明训练窗口的信息不足会限制机理扩展的收益。van Sluijs 等利用最优实验设计增加酶反应网络的参数信息量 \cite{vansluijs2024}，为后续 TMP 实验提供了可借鉴的方向。针对选择性与失活的实验设计，可分别增加早期 M--D 转化区间和后期减速区间的观测。

\subsection{机理模型与数据驱动校正的分工}
数据驱动方法在密集曲线内插中表现良好，说明局部预测并不总需要完整动力学参数解释。PPK-FIP-GP 在当前比较中获得最低合并内插误差，但 ORIG+GP 的表现接近，表明残差校正可以补偿不同机理模型在训练窗口内的偏差。高斯过程偏差项不能自动区分观测误差与模型缺项，因而残差校正后的高精度不构成机理正确性的独立证据。

前半段训练的外推中，PPK-FIP 优于本文设置的四种数据驱动基线。结论适用于有限样本、当前输入变量及相应训练方案。TabPFN、TabICL 和混合神经微分方程的近期研究表明，数据驱动方法的归纳约束与训练策略仍在发展 \cite{hollmann2025,qu2025tabicl,giampiccolo2024}。更明确的化学计量约束或反应速率参数化可能改变神经模型的外推表现，本文比较未覆盖所有物理约束学习方法。

\subsection{过程应用与研究局限}
PPK-FIP 可用于比较 TMP 酯化条件下的组成路径和动力学预测，但参数迁移需要保留酶--底物体系与反应条件的可比性。Tao 的跨水活度结果支持共享部分组成路径，Åkerman 的池化结果则提示温度变化可能限制参数共享。生物润滑剂过程设计应同时报告选择性估计、验证窗口和预测误差，并在新反应器条件下重新检验传质相关参数。

数据来源、验证规模和模型假设限制了结论范围。全部曲线来自已发表研究的数字化数据，缺少原始重复实验与独立误差估计，尤其无法确定 Mao 第一批次后期平台偏差的来源。四项研究形成六个拟合案例，留一案例超参数选择可能仍使用同一研究的其他条件，不能替代完全独立研究上的外部验证。边界参数的线性化置信区间也可能偏离名义覆盖率。完整参数集的实际可辨识性、先验宽度的跨体系适用性以及混合模型的守恒性质均需进一步验证。现有证据支持 PPK-FIP 在给定公开曲线上的预测改善，并支持将伪持续性和批内失活作为可检验的有效机理。
